% EM structure and teaching revision 20260928
\subsection{从实验数据判断模型范围}
粒子垂直磁场运动时，测得曲率半径直接给出 $p=|q|Br$。低速区用 $K=p^2/(2m)$，而相对论关系为
\[
 K=\sqrt{m^2c^4+p^2c^2}-mc^2.
\]
把根式按 $p/(mc)$ 展开，首项才是经典动能。因此“半径公式仍成立”不代表经典能量和周期公式仍成立；高能粒子的回旋角频率为 $|q|B/(\gamma m)$。
若粒子由电压差大小 $U$ 加速且初始静止，经典模型先给出 $K=|q|U$，再检查 $|q|U\ll mc^2$。对于电子，静能约为 $511\,\mathrm{keV}$：几千电子伏特的加速通常可作低速近似，而达到数百千电子伏特时不能再忽略相对论修正。
有界磁场问题还要检查几何可达性。先确定圆心、半径和真实绕行方向，再比较候选交点的到达时间；切线接触边界与穿越边界不同。若粒子在到达某个代数交点之前已离开磁区，后续轨迹就必须换用新区的运动方程。

